\documentclass[11pt]{article}

\usepackage[T1]{fontenc}
\usepackage[margin=1in]{geometry}
\usepackage{xcolor}
\usepackage{listings}
\usepackage{hyperref}

\lstset{
  basicstyle=\ttfamily\small,
  breaklines=true,
  columns=fullflexible,
  keepspaces=true,
  upquote=true,
  frame=single,
  showstringspaces=false
}

\title{Code and Data Guide\\\large modal-depth-experiment}
\author{}
\date{}

\begin{document}
\maketitle

\section*{Scope and reading guide}

This guide explains the seven top-level Haskell source files and the
\texttt{ScenarioInputs.csv} table. Haskell is reproduced in functional
sections, with an explanation directly beneath each excerpt. The CSV has
1,000,000 data rows; rather than repeat a very large, highly repetitive table,
its schema, row groups, and representative boundary rows are shown here.
Generated compiler files and build output are not project source files and are
not included.

\subsection{\texttt{Main.hs}}

\subsubsection{Build notes, module, imports, and modal profiles}

\begin{lstlisting}
-- ghc -package process -package directory -package filepath -package time -O2 -o sim_test Main.hs && ./sim_test
-- gzip -dc scenario-logs/ "the file in which you want to open up here" | less

-- highDangerValues = [7.125, 7.750, 8.125, 8.625, 9.000, 9.375, 9.625, 9.875, 7.375, 7.875, 8.375, 8.875, 9.125, 9.500, 9.750, 9.250]
-- lowDangerValues = [0.125, 0.375, 0.625, 0.875, 1.125, 1.375, 1.625, 1.875, 2.125, 2.375, 2.625, 2.875, 3.125, 3.375, 3.625, 3.875]

module Main where

import Types (ModeSet(..), Action(..), Result(uncertain))
import Experiment (runSimulation, runSafetyMethodSpeed)
import Metrics (countScored, accuracy, accuracyForExpected, accuracyForUncertain, accuracyForSpecification, countExpected, countFP, countFN, countUncertainAct, countUncertainNoAct)

{- Setting each modal depth through modeset -}
unimodal = ModeSet True False False False
bimodal  = ModeSet True True False False
trimodal = ModeSet True True True False
quadmodal = ModeSet True True True True
\end{lstlisting}

The first comments provide a direct GHC build/run command, a hint for
inspecting compressed trace output, and historical threshold lists. They do
not execute. The module imports the simulation, result types, and reporting
metrics. The four mode values progressively enable one through four
modalities; the current \texttt{main} selects \texttt{quadmodal}.

\subsubsection{Simulation and accuracy report}

\begin{lstlisting}
{-This section defines how the simulation is run-}
main :: IO ()
main = do
    (tracePath, results) <- runSimulation quadmodal 1000000

    putStrLn ("Safety trace archive: " ++ tracePath)
    print (length results)
    putStrLn "Conventional accuracy (all scenarios):"
    print (accuracy results)
    putStrLn "False positives:"
    print (countFP results)
    putStrLn "False negatives:"
    print (countFN results)
    putStrLn "Accuracy by expected outcome:"
    putStrLn "Expected ACT scenarios:"
    putStrLn "Scenario count:"
    print (countExpected ACT results)
    putStrLn "Accuracy:"
    print (accuracyForExpected ACT results)
    putStrLn "Expected NO_ACT scenarios:"
    putStrLn "Scenario count:"
    print (countExpected NO_ACT results)
    putStrLn "Accuracy:"
    print (accuracyForExpected NO_ACT results)
    putStrLn "Uncertain-zone scenarios:"
    putStrLn "Scenario count:"
    print (length (filter uncertain results))
    putStrLn "Accuracy:"
    print (accuracyForUncertain results)

    putStrLn "Specification Conformance (Did the Safety Method do what it was made to do?):"
    putStrLn "Scored scenarios:"
    print (countScored results)
    putStrLn "Accuracy:"
    print (accuracyForSpecification results)
    putStrLn "Uncertain-zone decisions:"
    putStrLn "ACT Scenarios:"
    print (countUncertainAct results)
    putStrLn "NO_ACT Scenarios:"
    print (countUncertainNoAct results)
    putStrLn "Expected: 100% specification conformance, conventional accuracy below 100%, and 340,000 ACT decisions in the uncertain zone."
\end{lstlisting}

\texttt{main} runs one million scenarios and saves a compressed safety trace.
It then prints overall conventional accuracy, false-positive and
false-negative counts, accuracy split by expected action, and uncertain-zone
counts and accuracy. The specification-specific score excludes uncertain
cases, while the final counters show which actions the policy took inside
that zone. The final printed expectation is an informational benchmark target,
not a runtime assertion.

\subsubsection{Safety-method timing report}

\begin{lstlisting}
    (speedCount, trialCount, meanNs, plusMinusNs) <- runSafetyMethodSpeed
    putStrLn "Safety method speed (CSV scenarios, no accuracy scoring):"
    putStrLn "Decision cycles per trial:"
    print speedCount
    putStrLn "Trials:"
    putStrLn ("N = " ++ show trialCount)
    putStrLn "Num Decision Cycles per Trial:"
    print(length results)
    putStrLn "Total Decision Cycles Count:"
    print(trialCount * length results)
    putStrLn "Average time per decision cycle:"
    putStrLn (show (round meanNs) ++ " ms ± " ++ show (round plusMinusNs) ++ " ms of cpu mutator time")
\end{lstlisting}

After the accuracy report, the program times the policy over CSV scenarios
without scoring accuracy. It prints the number of scenarios in a trial, the
trial count, total decision cycles, and the mean per-cycle time with a
half-range across trials. The returned time values are nanoseconds, although
the displayed labels say milliseconds; the displayed unit text should
therefore be interpreted cautiously.

\subsection{\texttt{Types.hs}}

\subsubsection{Language extension and scenario/mode/action inputs}

\begin{lstlisting}
{-# LANGUAGE UnboxedTuples #-}

module Types where

-- The UNPACK pragma removes the heap pointer and stores the raw Double directly inside the record.
-- The '!' ensures the value is evaluated instantly, preventing lazy thunks.
data Scenario = Scenario
  { c                :: {-# UNPACK #-} !Double
  , d_a              :: {-# UNPACK #-} !Double
  , d_e              :: {-# UNPACK #-} !Double
  , d_c              :: {-# UNPACK #-} !Double
  , d_s              :: {-# UNPACK #-} !Double
  , observed_context :: {-# UNPACK #-} !Double
  } deriving (Show, Eq)

data ModeSet = ModeSet
  { d_a_active       :: !Bool
  , d_e_active       :: !Bool
  , d_c_active       :: !Bool
  , d_s_active       :: !Bool
  } deriving (Show, Eq)

data Action = ACT | NO_ACT deriving (Show, Eq)
\end{lstlisting}

The extension enables unboxed-tuple syntax (none is declared in this module).
\texttt{Scenario} carries context and four modality values, with strict,
unpacked numeric fields to reduce lazy allocation overhead. \texttt{ModeSet}
records which modalities are enabled. \texttt{Action} is the binary decision
type shared throughout simulation and reporting.

\subsubsection{Thresholds and policy classification state}

\begin{lstlisting}
data SpecificThreshold = SpecificThreshold
  { highValues :: [Double]
  , lowValues :: [Double]
  } deriving (Show, Eq)

data DangerState = LowDanger | UncertainZone | HighDanger deriving (Show, Eq)

data PolicyState = PolicyState
  { stateA :: DangerState
  , stateE :: DangerState
  , stateC :: DangerState
  , stateS :: DangerState
  } deriving (Show, Eq)
\end{lstlisting}

\texttt{SpecificThreshold} holds the discrete values classified as high or
low danger. A value not in either list can be represented as
\texttt{UncertainZone}. \texttt{PolicyState} retains one danger label for
each modality.

\subsubsection{Trace and scored result records}

\begin{lstlisting}
data ProbabilityTrace = ProbabilityTrace
  { beforeProbabilities :: ![Double]
  , afterProbabilities :: ![Double]
  , noOpProbability :: !Double
  } deriving (Show, Eq)

data Result = Result
  { action           :: !Action
  , expectedAction   :: !Action
  , correct          :: !Bool
  , uncertain        :: !Bool
  } deriving (Show, Eq)
\end{lstlisting}

\texttt{ProbabilityTrace} stores model probabilities before and after the
safety intervention, plus a separate no-op probability. \texttt{Result}
keeps the actual and expected actions, whether they agree, and whether the
scenario belongs to the uncertain region.

\subsection{\texttt{Policy.hs}}

\subsubsection{Module imports and configured threshold lists}

\begin{lstlisting}
{- HLINT ignore "Use camelCase" -}
module Policy where

import Data.List (sort, sortOn)
import Types (Scenario(d_a, d_e, d_c, d_s), Action(..), SpecificThreshold(..), DangerState(..), PolicyState(..), ProbabilityTrace(..))

specificThreshold :: SpecificThreshold
specificThreshold = SpecificThreshold
  { highValues = [7.125, 7.750, 8.125, 8.625, 9.000, 9.375, 9.625, 9.875, 7.375, 7.875, 8.375, 8.875, 9.125, 9.500, 9.750, 9.250]
  , lowValues = [0.125, 0.375, 0.625, 0.875, 1.125, 1.375, 1.625, 1.875, 2.125, 2.375, 2.625, 2.875, 3.125, 3.375, 3.625, 3.875]
  }
\end{lstlisting}

The HLINT directive suppresses a naming-style suggestion for this module.
The policy imports sorting functions and the types it consumes or produces.
\texttt{specificThreshold} defines the exact discrete sets used to categorize
each modality.

\subsubsection{Danger classification and state refresh}

\begin{lstlisting}
classify :: SpecificThreshold -> Double -> DangerState
classify threshold value
  | value `elem` highValues threshold = HighDanger
  | value `elem` lowValues threshold = LowDanger
  | otherwise = UncertainZone

classifyForSafety :: SpecificThreshold -> Double -> DangerState
classifyForSafety threshold value =
  case classify threshold value of
    UncertainZone -> HighDanger
    danger -> danger

initialPolicyState :: PolicyState
initialPolicyState = PolicyState UncertainZone UncertainZone UncertainZone UncertainZone

updateStates :: PolicyState -> Scenario -> PolicyState
updateStates _ scenario = PolicyState
  { stateA = classifyForSafety specificThreshold (d_a scenario)
  , stateE = classifyForSafety specificThreshold (d_e scenario)
  , stateC = classifyForSafety specificThreshold (d_c scenario)
  , stateS = classifyForSafety specificThreshold (d_s scenario)
  }
\end{lstlisting}

\texttt{classify} marks exact high-list values as dangerous, exact low-list
values as safe, and all other values uncertain. The safety-specific wrapper
treats uncertainty as high danger (fail-safe). The initial state is uncertain
for every modality; each call to \texttt{updateStates} replaces all four
labels from the current scenario, so its prior-state argument is unused.

\subsubsection{Probability conversion and index helpers}

\begin{lstlisting}
softmax :: [Double] -> [Double]
softmax logits =
  let maximumLogit = maximum logits
      weights = map (exp . subtract maximumLogit) logits
      totalWeight = sum weights
  in map (\weight -> weight * 100.0 / totalWeight) weights

maximumIndex :: [(Int, Double)] -> Int
maximumIndex [] = 0
maximumIndex ((firstIndex, firstValue):remaining) =
  fst (foldl choose (firstIndex, firstValue) remaining)
  where
    choose best@(_, bestValue) candidate@(_, candidateValue)
      | candidateValue > bestValue = candidate
      | otherwise = best

normalize :: [Double] -> [Double]
normalize values =
  let total = sum values
  in map (\value -> value * 100.0 / total) values

assignProbabilities :: Int -> [(Int, Double)] -> [Double]
assignProbabilities count assignments =
  [ maybe 0.0 id (lookup index assignments)
  | index <- [0 .. count - 1]
  ]
\end{lstlisting}

\texttt{softmax} turns logits into percentages, subtracting the maximum first
to improve numerical stability. \texttt{maximumIndex} finds the index of the
largest value (and uses index zero for an empty input). \texttt{normalize}
rescales a list to sum to 100. \texttt{assignProbabilities} builds an
index-ordered list from assignments, using zero for missing indices.

\subsubsection{Rerouting probability mass}

\begin{lstlisting}
rerouteProbabilities :: [Double] -> Int -> Int -> [Double]
rerouteProbabilities probabilities dangerousIndex safeIndex =
  let orderedProbabilities = sort probabilities
      middleProbabilities = take (length probabilities - 2) (drop 1 orderedProbabilities)
      remainingIndices = sortOn (probabilities !!)
        [index | index <- [0 .. length probabilities - 1], index /= dangerousIndex, index /= safeIndex]
      assignments = (dangerousIndex, head orderedProbabilities)
        : (safeIndex, last orderedProbabilities)
        : zip remainingIndices middleProbabilities
  in normalize (assignProbabilities (length probabilities) assignments)
\end{lstlisting}

This helper sorts the probability values and assigns the smallest probability
to the selected dangerous index and the largest to the selected safe index.
Intermediate probabilities are paired with the other indices in their
original probability order, then the result is normalized. It assumes there
are enough probabilities for \texttt{head}, \texttt{last}, and the two
designated indices.

\subsubsection{No-op probability injection}

\begin{lstlisting}
fixedNoOpLogit :: Double
fixedNoOpLogit = case lowValues specificThreshold of
  value:_ -> value
  [] -> 0.0

injectNoOpProbabilities :: [Double] -> [Double]
injectNoOpProbabilities logits =
  let injectedProbabilities = softmax (logits ++ [fixedNoOpLogit])
      orderedProbabilities = sort injectedProbabilities
      existingCount = length logits
      existingIndices = sortOn (injectedProbabilities !!) [0 .. existingCount - 1]
      assignments = zip existingIndices (take existingCount orderedProbabilities)
        ++ [(existingCount, last orderedProbabilities)]
  in normalize (assignProbabilities (existingCount + 1) assignments)
\end{lstlisting}

The no-op option gets a fixed logit equal to the first low-danger threshold.
After calculating softmax across the original actions and this extra option,
the probabilities are reassigned in sorted order: the no-op receives the
largest probability, while the original action slots receive the remaining
values ordered against their original probabilities. The result is
renormalized.

\subsubsection{Safety decision and top-level policy}

\begin{lstlisting}
execute :: PolicyState -> [Double] -> (Action, [Double], [Double], Double)
execute states logits =
  let labels = [stateA states, stateE states, stateC states, stateS states]
      probabilities = softmax logits
      winner = maximumIndex (zip [0..] probabilities)
      safeIndices =
        [ (index, probability)
        | (index, (label, probability)) <- zip [0..] (zip labels probabilities)
        , label == LowDanger
        ]
  in if labels !! winner == LowDanger
       then (NO_ACT, probabilities, probabilities, 0.0)
       else case safeIndices of
         [] ->
           let injectedProbabilities = injectNoOpProbabilities logits
           in (ACT, probabilities, take (length logits) injectedProbabilities, last injectedProbabilities)
         _ ->
           let safestIndex = maximumIndex safeIndices
           in (ACT, probabilities, rerouteProbabilities probabilities winner safestIndex, 0.0)

policy :: PolicyState -> Scenario -> (PolicyState, Action, ProbabilityTrace)
policy previousState scenario =
  let nextState = updateStates previousState scenario
      logits = [d_a scenario, d_e scenario, d_c scenario, d_s scenario]
      (action, before, after, noOp) = execute nextState logits
      trace = ProbabilityTrace before after noOp
  in (sum before + sum after + noOp) `seq` (nextState, action, trace)
\end{lstlisting}

\texttt{execute} computes probabilities over the four modality logits. If the
highest-probability modality is safe, it returns \texttt{NO\_ACT} with
probabilities unchanged. Otherwise it returns \texttt{ACT}: if no modality is
safe, it adds the no-op probability; if a safe modality exists, it reroutes
probability mass toward the safest safe index and away from the winner.
\texttt{policy} refreshes classifications, extracts logits, calls
\texttt{execute}, creates a trace, and uses \texttt{seq} to force the sums
before returning the tuple.

\subsection{\texttt{Metrics.hs}}

\subsubsection{Accuracy calculations}

\begin{lstlisting}
module Metrics where

import Types (Result(..), Action(..))
import Evaluate (falsePositive, falseNegative)

accuracy :: [Result] -> Double
accuracy = accuracyWhere (const True)

accuracyForExpected :: Action -> [Result] -> Double
accuracyForExpected expected = accuracyWhere matchesExpected
    where
        matchesExpected result = not (uncertain result) && expectedAction result == expected

accuracyForUncertain :: [Result] -> Double
accuracyForUncertain = accuracyWhere uncertain

-- First 66% only: expected ACT and NO_ACT. Uncertain-zone scenarios are not scored.
accuracyForSpecification :: [Result] -> Double
accuracyForSpecification = accuracyWhere (not . uncertain)

accuracyWhere :: (Result -> Bool) -> [Result] -> Double
accuracyWhere include results =
    let selected = filter include results
        total = length selected
        correctCount = length (filter correct selected)
    in if total == 0 then 0.0 else fromIntegral correctCount / fromIntegral total
\end{lstlisting}

\texttt{accuracy} scores every result. \texttt{accuracyForExpected} selects
only certain scenarios with a specified expected action;
\texttt{accuracyForUncertain} selects only uncertain cases; and
\texttt{accuracyForSpecification} excludes uncertain cases. Their common
helper computes the fraction marked correct and returns zero if nothing was
selected.

\subsubsection{Counts and error categories}

\begin{lstlisting}
countScored :: [Result] -> Int
countScored = length . filter (not . uncertain)

countFP :: [Result] -> Int
countFP = length . filter falsePositive

countFN :: [Result] -> Int
countFN = length . filter falseNegative

countUncertainAct :: [Result] -> Int
countUncertainAct = length . filter isUncertainAct
    where
        isUncertainAct result = uncertain result && action result == ACT

countUncertainNoAct :: [Result] -> Int
countUncertainNoAct = length . filter isUncertainNoAct
    where
        isUncertainNoAct result = uncertain result && action result == NO_ACT

countExpected :: Action -> [Result] -> Int
countExpected expected = length . filter matchesExpected
    where
        matchesExpected result = not (uncertain result) && expectedAction result == expected
\end{lstlisting}

These functions count certain/scored scenarios, false positives, false
negatives, and each action taken in the uncertain zone. \texttt{countExpected}
counts only certain cases whose expected action matches its argument.

\subsection{\texttt{Evaluate.hs}}

\subsubsection{Scenario evaluation and error predicates}

\begin{lstlisting}
module Evaluate where

import Types ( Action(..), Scenario, Result(..), PolicyState, ProbabilityTrace )
import Policy (policy)

processScenario :: PolicyState -> (Scenario, Action, Bool) -> (PolicyState, Result, ProbabilityTrace)
processScenario previousState (scenario, expected, isUncertain) =
    let (nextState, actualAction, trace) = policy previousState scenario
    in (nextState, Result actualAction expected (actualAction == expected) isUncertain, trace)

falsePositive :: Result -> Bool
falsePositive r = action r == ACT && not (correct r)

falseNegative :: Result -> Bool
falseNegative r = action r == NO_ACT && not (correct r)
\end{lstlisting}

\texttt{processScenario} applies the policy to one input, then combines its
actual action with the expected action and uncertain flag to produce a
\texttt{Result}. A false positive is an incorrect \texttt{ACT}; a false
negative is an incorrect \texttt{NO\_ACT}.

\subsection{\texttt{Experiment.hs}}

\subsubsection{Running the simulation and streaming trace output}

\begin{lstlisting}
module Experiment where

import Control.Exception (evaluate, onException)
import Control.Monad (foldM, replicateM)
import Data.List (intercalate, foldl')
import Data.Time (defaultTimeLocale, formatTime, getCurrentTime)
import System.Directory (createDirectoryIfMissing, removeFile)
import System.CPUTime (getCPUTime)
import System.Exit (ExitCode(..))
import System.FilePath ((</>))
import System.IO (Handle, IOMode(WriteMode), hClose, hPutStr, withBinaryFile)
import System.IO.Error (catchIOError)
import System.Process (CreateProcess(..), StdStream(CreatePipe, UseHandle), proc, waitForProcess, withCreateProcess)
import Text.Printf (printf)
import Types (ModeSet, PolicyState(..), Result(..), Scenario(..), Action, ProbabilityTrace(..))
import Scenario (generateScenarios, loadScenarioCases)
import Evaluate (processScenario)
import Policy (initialPolicyState, policy)

runSimulation :: ModeSet -> Int -> IO (FilePath, [Result])
runSimulation _ n = do
    scenarios <- generateScenarios n
    createDirectoryIfMissing True "scenario-logs"
    timestamp <- formatTime defaultTimeLocale "%Y%m%d-%H%M%S%q" <$> getCurrentTime
    let tracePath = "scenario-logs" </> ("safety-trace-" ++ timestamp ++ ".txt.gz")
    results <- writeCompressedTrace tracePath $ \traceHandle -> do
        (_, reversedResults) <- foldM (processAndTrace traceHandle) (initialPolicyState, []) (zip [0..] scenarios)
        return (reverse reversedResults)
    return (tracePath, results)

processAndTrace
    :: Handle
    -> (PolicyState, [Result])
    -> (Int, (Scenario, Action, Bool))
    -> IO (PolicyState, [Result])
processAndTrace traceHandle (previousState, results) (index, scenarioCase@(scenario, _, _)) = do
    let (nextState, result, trace) = processScenario previousState scenarioCase
    hPutStr traceHandle (renderScenarioTrace index scenario result nextState trace)
    return (nextState, result : results)
\end{lstlisting}

\texttt{runSimulation} loads up to the requested number of cases, creates the
trace directory, builds a timestamped archive path, processes every scenario
in order, and returns the archive path and results. The mode argument is
currently ignored (as indicated by \texttt{\_}). \texttt{processAndTrace}
evaluates one case, appends its trace text, and carries forward policy state
and accumulated results.

\subsubsection{Formatting per-scenario traces}

\begin{lstlisting}
renderScenarioTrace :: Int -> Scenario -> Result -> PolicyState -> ProbabilityTrace -> String
renderScenarioTrace index scenario result states trace =
    intercalate "\n"
        [ "Scenario " ++ printf "%06d" (index + 1) ++ " (CSV row " ++ show (index + 2) ++ ")"
        , "Raw logits:"
        , renderValues [("d_a", d_a scenario), ("d_e", d_e scenario), ("d_c", d_c scenario), ("d_s", d_s scenario)]
        , "Danger classifications:"
        , renderValues [("d_a", stateA states), ("d_e", stateE states), ("d_c", stateC states), ("d_s", stateS states)]
        , "Expected action: " ++ show (expectedAction result)
        , "Uncertain zone: " ++ show (uncertain result)
        , "Probabilities before intervention:"
        , renderPercentages (beforeProbabilities trace)
        , "Safety decision: " ++ show (action result)
        , "Probabilities after intervention:"
        , renderPercentages (afterProbabilities trace)
        , "  no_op: " ++ formatPercentage (noOpProbability trace)
        , "  total: " ++ formatPercentage (sum (afterProbabilities trace) + noOpProbability trace)
        , ""
        ]

renderValues :: Show value => [(String, value)] -> String
renderValues = unlines . map (\(name, value) -> "  " ++ name ++ ": " ++ show value)

renderPercentages :: [Double] -> String
renderPercentages = unlines . map (\(name, value) -> "  " ++ name ++ ": " ++ formatPercentage value)
    . zip ["d_a", "d_e", "d_c", "d_s"]

formatPercentage :: Double -> String
formatPercentage = printf "%.6f%%"
\end{lstlisting}

\texttt{renderScenarioTrace} emits a readable record with a one-based scenario
number, corresponding CSV row, logits, classifications, expected and actual
decisions, uncertainty status, and probability values before and after
intervention. \texttt{renderValues} formats named values;
\texttt{renderPercentages} associates the four probability positions with
their modality names; and \texttt{formatPercentage} prints six fractional
digits and a percent sign.

\subsubsection{Gzip archive lifecycle and failure handling}

\begin{lstlisting}
writeCompressedTrace :: FilePath -> (Handle -> IO result) -> IO result
writeCompressedTrace tracePath writeTrace =
    writeArchive `onException` (removeFile tracePath `catchIOError` (const (return ())))
  where
    writeArchive = withBinaryFile tracePath WriteMode $ \archiveHandle ->
        withCreateProcess
            ((proc "gzip" ["-c"]) { std_in = CreatePipe, std_out = UseHandle archiveHandle })
            (\inputHandle _ _ processHandle -> do
                traceHandle <- maybe (ioError (userError "gzip did not provide an input stream")) return inputHandle
                result <- writeTrace traceHandle
                hClose traceHandle
                exitCode <- waitForProcess processHandle
                case exitCode of
                    ExitSuccess -> return result
                    ExitFailure code -> ioError (userError ("gzip exited with status " ++ show code)))
\end{lstlisting}

This function pipes trace text into a \texttt{gzip -c} child process while
writing its compressed output to the archive file. It closes the pipe,
checks gzip's exit status, and raises an IO error if gzip fails or does not
provide stdin. If any operation throws, it attempts to remove the incomplete
archive; an error from that cleanup attempt is ignored.

\subsubsection{Repeated policy-only timing}

\begin{lstlisting}
speedTrialCount :: Int
speedTrialCount = 100

-- Time only the safety method (policy) over the CSV scenarios. Accuracy
-- scoring is skipped so the result is per-decision-cycle runtime.
-- Returns (cycles per trial, trial count, mean ns/cycle, half-range ns).
runSafetyMethodSpeed :: IO (Int, Int, Double, Double)
runSafetyMethodSpeed = do
    cases <- loadScenarioCases
    let scenarios = map fst cases
    _ <- evaluate (foldl' (\() scenario -> scenario `seq` ()) () scenarios)
    perCycleNanos <- replicateM speedTrialCount (timePerCycleNanos scenarios)
    let meanNs = sum perCycleNanos / fromIntegral speedTrialCount
        plusMinusNs = (maximum perCycleNanos - minimum perCycleNanos) / 2.0
    return (length scenarios, speedTrialCount, meanNs, plusMinusNs)

timePerCycleNanos :: [Scenario] -> IO Double
timePerCycleNanos scenarios = do
    start <- getCPUTime
    let finalState = foldl' applyPolicy initialPolicyState scenarios
    _ <- evaluate (forcePolicyState finalState)
    end <- getCPUTime
    let totalSeconds = fromInteger (end - start) / 1.0e12
    return (totalSeconds / fromIntegral (length scenarios) * 1.0e9)

applyPolicy :: PolicyState -> Scenario -> PolicyState
applyPolicy previousState scenario =
    let (nextState, action, trace) = policy previousState scenario
    in action `seq` trace `seq` forcePolicyState nextState

forcePolicyState :: PolicyState -> PolicyState
forcePolicyState states =
    stateA states `seq`
    stateE states `seq`
    stateC states `seq`
    stateS states `seq`
    states
\end{lstlisting}

\texttt{runSafetyMethodSpeed} loads all CSV cases, extracts scenarios, forces
them before timing, then performs 100 independent timing trials. Each trial
uses a strict fold through \texttt{policy}, forcing the resulting state,
action, and trace so that lazy work is included. \texttt{timePerCycleNanos}
measures CPU time and converts the average duration per scenario to
nanoseconds. The returned plus/minus figure is half the observed range, not a
statistical confidence interval.

\subsection{\texttt{Scenario.hs}}

\subsubsection{Configuration and CSV header}

\begin{lstlisting}
module Scenario where
import Types (Scenario(..), Action(..), SpecificThreshold(..))
import Text.Read (readMaybe)
import Policy (specificThreshold)
import System.IO (IOMode(WriteMode), hPutStrLn, withFile)
import Text.Printf (printf)

scenarioContext :: Double
scenarioContext = 0.1

scenarioCount :: Int
scenarioCount = 1000000

unidentifiedScenarioValues :: [Double]
unidentifiedScenarioValues = [4.125, 4.500, 4.875, 5.250, 5.625, 6.000, 6.375, 6.750]

scenarioCsvHeader :: String
scenarioCsvHeader = "d_a,d_e,d_c,d_s,expected_action"
\end{lstlisting}

These constants define the context value attached to loaded scenarios, the
expected number of table rows, and the intermediate modality values used in
the uncertain portion. The header fixes the CSV column order: four modality
values followed by the expected action.

\subsubsection{Loading and validating the CSV}

\begin{lstlisting}
loadScenarioCases :: IO [(Scenario, Action)]
loadScenarioCases = do
    contents <- readFile "ScenarioInputs.csv"
    case lines contents of
        header:rows
            | header == scenarioCsvHeader -> do
                cases <- either (ioError . userError) return
                    (traverse (uncurry parseScenarioRow) (zip [2..] rows))
                if length cases == scenarioCount
                    then return cases
                    else ioError (userError ("Expected " ++ show scenarioCount ++ " scenario rows, found " ++ show (length cases)))
            | otherwise -> ioError (userError "ScenarioInputs.csv has an unexpected header")
        [] -> ioError (userError "ScenarioInputs.csv is empty")

parseScenarioRow :: Int -> String -> Either String (Scenario, Action)
parseScenarioRow rowNumber row =
    case splitCsvRow row of
        [valueA, valueE, valueC, valueS, expectedText] -> do
            modalityA <- parseModality "d_a" valueA
            modalityE <- parseModality "d_e" valueE
            modalityC <- parseModality "d_c" valueC
            modalityS <- parseModality "d_s" valueS
            expected <- case expectedText of
                "ACT" -> Right ACT
                "NO_ACT" -> Right NO_ACT
                _ -> Left (rowError "expected_action must be ACT or NO_ACT")
            return
                ( Scenario scenarioContext modalityA modalityE modalityC modalityS scenarioContext
                , expected
                )
        _ -> Left (rowError "expected five comma-separated fields")
  where
    rowError message = "ScenarioInputs.csv row " ++ show rowNumber ++ ": " ++ message
    parseModality name valueText =
        case (readMaybe valueText :: Maybe Double) of
            Just value
                | not (isNaN value || isInfinite value) && value >= 0.0 && value <= 10.0 -> Right value
            _ -> Left (rowError (name ++ " must be a number between 0 and 10"))

splitCsvRow :: String -> [String]
splitCsvRow row =
    case break (== ',') row of
        (field, []) -> [field]
        (field, _:remaining) -> field : splitCsvRow remaining
\end{lstlisting}

\texttt{loadScenarioCases} reads the file, verifies the exact header, parses
every row with its physical CSV line number, and rejects any row count other
than one million. \texttt{parseScenarioRow} requires exactly five comma
separated values, converts each modality to a finite number in \([0,10]\),
checks the action spelling, and constructs a \texttt{Scenario}. Parse
problems include the CSV row number. The small \texttt{splitCsvRow} helper
splits on commas; it does not implement quoted-field CSV escaping.

\subsubsection{Marking uncertain cases}

\begin{lstlisting}
generateScenarios :: Int -> IO [(Scenario, Action, Bool)]
generateScenarios n = do
    cases <- loadScenarioCases
    let count = min scenarioCount (max 0 n)
    return
        [ (scenario, expected, index >= 660000)
        | (index, (scenario, expected)) <- zip [0..] (take count cases)
        ]
\end{lstlisting}

The requested count is clamped to the range zero through one million.
Each output item combines its parsed scenario and expected action with an
uncertainty flag, which becomes true beginning at zero-based index 660,000
(CSV row 660,002).

\subsubsection{Writing and composing the generated table}

\begin{lstlisting}
writeScenarioInputs :: FilePath -> IO ()
writeScenarioInputs path = withFile path WriteMode $ \handle -> do
    hPutStrLn handle scenarioCsvHeader
    mapM_ (hPutStrLn handle . renderScenarioRow) [0 .. scenarioCount - 1]

renderScenarioRow :: Int -> String
renderScenarioRow index =
    let threshold = specificThreshold
        values
          | index < 330000 =
              insertAt (index `mod` 4)
                (highValues threshold !! (index `mod` length (highValues threshold)))
                (takeRotated 3 (lowValues threshold) (index * 3))
          | index < 660000 = takeRotated 4 (lowValues threshold) (index * 4)
          | index < 830000 =
              insertAt ((index - 660000) `mod` 4)
                (unidentifiedScenarioValues !! ((index - 660000) `mod` length unidentifiedScenarioValues))
                (takeRotated 3 (lowValues threshold) ((index - 660000) * 3))
          | otherwise = takeRotated 4 unidentifiedScenarioValues ((index - 830000) * 4)
        expected
          | index < 330000 = "ACT"
          | index < 660000 = "NO_ACT"
          | index `mod` 2 == 0 = "ACT"
          | otherwise = "NO_ACT"
    in joinWithComma (map (printf "%.3f") values ++ [expected])

takeRotated :: Int -> [a] -> Int -> [a]
takeRotated count values offset = take count (drop (offset `mod` length values) (cycle values))

insertAt :: Int -> a -> [a] -> [a]
insertAt index value values =
    take index values ++ [value] ++ drop index values

joinWithComma :: [String] -> String
joinWithComma [] = ""
joinWithComma [value] = value
joinWithComma (value:remaining) = value ++ "," ++ joinWithComma remaining
\end{lstlisting}

\texttt{writeScenarioInputs} writes the header and one generated record for
each index from zero through 999,999. \texttt{renderScenarioRow} defines four
contiguous blocks: one high plus three low values, four low values, one
unidentified plus three low values, and four unidentified values. Expected
actions are always \texttt{ACT} in the first block, always \texttt{NO\_ACT}
in the second, and alternate in the last two. Values are formatted to three
decimal places. \texttt{takeRotated} cycles through a list from an offset,
\texttt{insertAt} inserts the designated value, and \texttt{joinWithComma}
serializes the fields.

\subsection{\texttt{ScenarioInputs.csv}}

\subsubsection{Header and record shape}

\begin{lstlisting}[basicstyle=\ttfamily\small]
d_a,d_e,d_c,d_s,expected_action
7.125,0.125,0.375,0.625,ACT
0.875,7.750,1.125,1.375,ACT
1.625,1.875,8.125,2.125,ACT
2.375,2.625,2.875,8.625,ACT
\end{lstlisting}

The first line names four modality columns and the ground-truth
\texttt{expected\_action}. Each subsequent line has four numeric modality
values and one \texttt{ACT} or \texttt{NO\_ACT} label. The displayed records
are the table's first four data rows (CSV rows 2--5).

\subsubsection{Four contiguous data groups}

\begin{lstlisting}[basicstyle=\ttfamily\small]
CSV rows 2--330001:
  one high-threshold value and three low-threshold values; expected ACT

CSV rows 330002--660001:
  four low-threshold values; expected NO_ACT

CSV rows 660002--830001:
  one unidentified/intermediate value and three low-threshold values;
  expected actions alternate ACT, NO_ACT

CSV rows 830002--1000001:
  four unidentified/intermediate values;
  expected actions alternate ACT, NO_ACT
\end{lstlisting}

These row ranges include the header's offset: the first data record is CSV
row 2. The ranges contain 330,000, 330,000, 170,000, and 170,000 records,
respectively, totaling 1,000,000. In the supplied file, there are 500,000
\texttt{ACT} labels and 500,000 \texttt{NO\_ACT} labels. Values are drawn
from repeating combinations; the file contains 38 distinct four-value
modality tuples, with many rows repeating those tuples.

\subsubsection{Rows at group boundaries and end of file}

\begin{lstlisting}[basicstyle=\ttfamily\small]
CSV row 330001:  3.375,3.625,3.875,9.250,ACT
CSV row 330002:  0.125,0.375,0.625,0.875,NO_ACT
CSV row 660001:  3.125,3.375,3.625,3.875,NO_ACT
CSV row 660002:  4.125,0.125,0.375,0.625,ACT
CSV row 830001:  3.375,3.625,3.875,6.750,NO_ACT
CSV row 830002:  4.125,4.500,4.875,5.250,ACT
CSV row 1000001: 5.625,6.000,6.375,6.750,NO_ACT
\end{lstlisting}

These examples show the transitions between the four blocks and the final
record. Intermediate values in the last two blocks are not members of the
configured low/high threshold lists, so the policy's safety classification
treats them as high danger even though the CSV marks the scenarios as part of
the uncertain region for separate scoring.

\end{document}
